\section{Logical Spine and Proof Ledger}
\label{sec:logical_spine}

\subsection{The Proof Ledger}
To ensure the auditability of the main claim (Existence of a Mass Gap), we present the logical dependency graph as a formal "Proof Ledger". This structure allows referees to discharge verify assumptions one by one.

\begin{framed}
\textbf{Theorem 1 (Main Result):} 4D Yang-Mills has a strictly positive mass gap $\Delta > 0$.
\end{framed}

\begin{longtable}{|p{0.3\textwidth}|p{0.4\textwidth}|p{0.2\textwidth}|}
\hline
\textbf{Component} & \textbf{Claim/Theorem} & \textbf{Method / Status} \\
\hline
\textbf{1. Strong Coupling Base} & Uniqueness of Mass Gap for $\beta \le \beta_{Dob}$. & \textbf{Analytic} (Dobrushin Criterion, App G). \newline \textit{Status: Proven for $\beta < 0.5$ (SU3).} \\
\hline
	extbf{2. Intermediate Bridge} & Contraction of RG Flow for $\beta \in [\VerBetaIntermediateMin, \VerBetaIntermediateMax]$. & \textbf{Computer-Assisted} (Interval Arithmetic). \newline \textit{Status: Verified (Sec 5.5).} \\
\hline
	extbf{3. Weak Coupling} & Uniform Log-Sobolev Inequality for $\beta > \VerBetaWeakMin$. & \textbf{Analytic} (Strict Convexity $\implies$ LSI). \newline \textit{Status: Validated (Sec 3.5).} \\
\hline
\textbf{4. Scaling Limit} & Mass Gap scales as $e^{-\beta/2\beta_0}$. & \textbf{Analytic} (Balaban/Symanzik). \newline \textit{Status: Standard QFT.} \\
\hline
\textbf{5. Lorentz Restoration} & Anisotropy $\xi \to 1$ fixed point exists. & \textbf{Hybrid} (Analytic Definition + CAP check of Jacobian $<1$). \newline \textit{Status: Verified (Thm 5.3).} \\
\hline
\end{longtable}

\subsection{Verification Roadmap}
\begin{figure}[h]
\centering
\caption{Theorem Dependency Graph (Visual)}
\end{figure}
$\uparrow$ \textit{(depends on)}
\begin{enumerate}
    \item \textbf{Uniform Log-Sobolev Inequality (LSI):} (Chapter 3)
         \begin{itemize}
             \item Proves $\Delta \ge \gamma_{LSI}$ uniform in volume.
             \item \textit{Method:} Analytic (Conditional Tensorization + Cluster Expansion).
         \end{itemize}
    \item \textbf{Regime Bridging:}
         \begin{itemize}
             \item \textbf{Weak Coupling:} Analytic (Cluster Expansion, Glimm-Jaffe-Spencer).
             \item \textbf{Strong Coupling:} Analytic (Standard High-T Expansion coverage).
             \item \textbf{Intermediate Regime:} Computer-Assisted Proof (CAP) -- see \S\ref{sec:cap_dependency}.
         \end{itemize}
    \item \textbf{Continuum Limit:} (Chapter 4)
         \begin{itemize}
             \item Norm Resolvent Convergence of Hamiltonians.
             \item \textit{Method:} Analytic (Symanzik analysis).
         \end{itemize}
\end{enumerate}

\subsection{The Role of the Computer-Assisted Proof}
\label{sec:cap_dependency}
The proof is \textbf{hybrid}. The analytic framework (LSI, RG flow) provides the architecture, but the \textit{connection} between the strong-coupling phase (proved by expansions) and the weak-coupling phase (proved by perturbation theory) relies on the CAP.

\begin{itemize}
    \item \textbf{What CAP Certifies:} It rigorously verifies the \textit{Tube Contraction Condition} (Theorem K.1, Appendix~\ref{app:interval_arithmetic}) for the intermediate effective action, ensuring no phase transitions or critical points occur in the crossover interaction range $g \in [g_{strong}, g_{weak}]$.
    \item \textbf{Rigor Level:} The CAP uses \textbf{Interval Arithmetic} with IEEE 754 directed rounding to produce mathematical certificates. It is not a numerical simulation. See Section~\ref{sec:cap_formal} for the formal specification.
    \item \textbf{Dependency:} The main theorem for SU(2)/SU(3) depends essentially on this certificate.
\end{itemize}

\begin{remark}[Terminology: "Rigorous Computer-Assisted" vs. "Unconditional"]
\label{rmk:unconditional}
We use "unconditional" to mean the proof does not rely on unproven mathematical conjectures (such as the Riemann Hypothesis or a generalization of existing theorems). However, the proof \textbf{does} rely on the correctness of:
\begin{enumerate}
    \item The CAP implementation and its certificate verification (auditable via \texttt{interval\_gap\_check.py}).
    \item The interval arithmetic library (mpmath) correctly implementing IEEE 754.
    \item The underlying hardware correctly executing IEEE 754 operations.
\end{enumerate}
This is standard for computer-assisted proofs (cf. Hales's proof of Kepler, Tucker's proof of Lorenz attractor). The certificate is \textbf{independently checkable} by running the minimal verifier on any IEEE 754-compliant platform.
\end{remark}

\subsection{Proof Component Summary}
\begin{longtable}{|p{0.25\textwidth}|p{0.4\textwidth}|p{0.25\textwidth}|}
\hline
\textbf{Component} & \textbf{Statement} & \textbf{Method} \\ \hline
\textbf{Lattice Action} & SU(2)/SU(3): Wilson Action \newline SU(N $\ge$ 5): Anisotropic Iwasaki-Wilson (App.~\ref{ch:action_definitions}) & Defined \\ \hline
\textbf{Strong Coupling} & Convergent Cluster Expansion & Analytic \\ \hline
\textbf{Weak Coupling} & Asymptotic Freedom + Renormalizable & Analytic \\ \hline
\textbf{Intermediate Bridge} & Tube Contraction (No Phase Transition) & CAP (Thm.~\ref{thm:cap_formal}) \\ \hline
\textbf{Uniform LSI} & Global Gap from Local Mixing & Analytic (Ch.~3) \\ \hline
\textbf{Gap Comparison} & $\Delta_H \ge c \cdot \gamma_{LSI}$ & Analytic (App.~\ref{app:gap_equivalence}) \\ \hline
\textbf{Continuum Limit} & Gap Stability under $a \to 0$ & Analytic (App.~D) \\ \hline
\end{longtable}
